\documentclass{article}
\usepackage{graphicx} % Required for inserting images
\usepackage[table]{xcolor} % For \rowcolor
\usepackage{siunitx}[=v2] % For 's' column type in tabular
\usepackage[utf8]{inputenc}
\usepackage[T1]{fontenc}
\usepackage{array}
\usepackage{tabularx}
\usepackage[table]{xcolor}
\usepackage{xcolor}
\usepackage[dvipsnames]{xcolor}

\title{Lastenheft}
\author{Nisa, Esra, Rebeca, Timon, Felix, Tobias, Elmo}
\date{November 2026}

\begin{document}

\maketitle
 

\section{Auftraggeber}
Auftraggeber: Steve Durstewitz M.sc, Prof. Dr. Claudia Lenk, Insitut BMT 


\section{Einführung}
\subsection{Ausgangszustand und Ist-Zustand}
Das Open-Earable ist ein Open-Source Projekt des KIT -Karlsruher Institute of Technology- und umfasst die Hardware mit zwei Kopfhörer und das Github-Projekt mit der passenden Software. Die bisher verfügbaren Hardware Komponenten umfassen einen . Die bisherigen Softwareimplementierungen umfassen . 
Die bisherigen Messwerte und Ergebnisse der Forschung des vorherigen Projektes sind uns zur Verfügung und sind die Basis unserer Aufgabenbereitung. 
Jedgliche Hardwarebestandteile werden vollständig übernommen und die Basis ist die Software des Herstellers.
Der aktuelle Leistungsstand der Geräuschunterdrückung und Sprachverarbeitung ist zu Projektbeginn nicht vollständig quantifiziert. Daher soll eine reproduzierbare messtechnische Untersuchung des bestehenden Systems als Grundlage für die weitere Softwareentwicklung erfolgen


\subsection{Motivation}
%Im Jahre 2025 beläuft sich die Anzahl der meldepflichtigen Arbeitsunfälle auf Rund 731 Tausend. Diese Zahl ist zwar im Vergleich zum Vorjahr gesunken, jedoch ist die Arbeit immer noch der dritthäufigste Ort an dem Unfälle geschehen. Jedoch beläuft sich diese Zahl nicht auf die subtilen Belastungen, die einem Arbeiter jedem Tag ausgesetzt sind, die sowohl körperliche als auch psychische (Folge-)Schäden auch längerfristig erhalten kann. Eine erhöhte Lärmbelästigung am Arbeitsplatz erhöht die Unfallwahrscheinlichkeit und ist für das Hörvermögen und psychische Unversehrtheit eine Gefahr. Vorallem Arbeitsumgebungen im Baugewerbe oder der Metallverarbeitung sind viele Umgebungsgeräusche Alltag der Arbeiter*Innen. Um die Gesundheit der Arbeitnehmer verstärkt zu Schützen wollen wir an der Filterung und Signalverarbeitung arbeiten. In diesem Zuge ist unser zweiter Fokus auf der Verarbeitung der Sprache zu lesbarem Text kommt auch zu Nutze.

Der Arbeitsplatz ist der dritthäufigste Ort an dem Unfälle entstehen. Im Jahr 2025 wurden in Deutschland rund 731.000 meldepflichtige Arbeitsunfälle registriert. Neben akuten Unfallereignissen stellen auch langfristige Belastungen am Arbeitsplatz eine Gefahr für die körperliche und psychische Gesundheit von Beschäftigten dar. 

Insbesondere in lärmintensiven Arbeitsumgebungen, beispielsweise im Baugewerbe oder in der Metallverarbeitung, können erhöhte Schallpegel das Hörvermögen beeinträchtigen, die Kommunikation erschweren und das Unfallrisiko erhöhen.

Das Projekt untersucht daher die Möglichkeiten einer gezielten Unterdrückung von störenden Umgebungsgeräuschen unter Verwendung der OpenEarable-Plattform. Ergänzend soll die digitale Transkription gesprochener Sprache die Kommunikation in geräuschintensiven Umgebungen unterstützen.

\subsection{Projektziel und Soll-Zustand}
%Den jetzigen Stand der Signalfilterung zur Erheben anhand eigensdurchgeführter standartisierten Messungen und anhand dieser und einer Marktforschung die Tonverarbeitung anzupassen. Die Messdaten der vorherigen Gruppe werden in Betracht gezogen und mit unseren erhobenen Daten verglichen. Am Ende unseres Projektes sollen die Störgeräusche weniger vernommen werden und in einem Gespräch (auch in lauter Umgebung) die gesprochenen Worte als Text auf einer ausgewählten Bildausgabe verschriftlich werden.  
Unser Projekt wird sich mit den drei Zielen auseinandersetzten:
\begin{enumerate}
    \item Ziel 1: Geräuschunterdrückung
    \subitem Reduktion störender akaustischer Signale, denen der Nutzer ausgesetzt ist
    \item Ziel 2: Sprachverständlichkeit
    \subitem Erhalt und Verbesserung der Verständlichkeit gesprochener Sprache bei Hintergrundgeräuschen
    \item Ziel 3: Sprachtranskription
    \subitem Umwandlung gesprocheneer Sprache in digital darstellbaren Text
\end{enumerate}

\subsection{Abgrenzung des Projekts}
% Ziel dieser Studierendenenarbeit ist es innerhalb der vorgegebenen Zeiten ein 
% zwei Untergruppen sich die folgenden Ziele zu erarbeiten: 
% \begin{itemize}
%     \item Erfassen alles theoretischen Grundlagen der technischen Umsetzung, inklusiver Marktforschung
%     \item Überprüfung der praktischen Umsetzung der Ziele im Rahmen der technisch- und softwarebegrenzenden Möglichkeiten
%     \item Umsetzung der Ziele in der Software
%     \item Durchführung Messungen zum Überprüfen der bisherigen Software und die Ergänzungen der Gruppe
% \end{itemize}

% Eine Untergruppe besteht aus . Diese beschäftigt sich mit Text-to-Speech, indem Gespräche in unmittelbarer Umgebung verarbeitet werden und in einer Textform ausgegeben wird. \\
% Die andere Untergruppe, bestehend aus Felix, Elmo, beschäftigt sich mit der Filterung des Umgebungslärmes, das im Arbeitumfeld zu einer großen Belastung und gesundheitlichen Bedrohung wird. Unser Ziel ist es den Hintergrundlärm zu messen, analysieren und eine Active Noise Cancelling und besser Geräuschfilterung zu implementieren.\\

% Dieses Lastenheft legt die Anforderungen für ein Gerät zur Schallmessung, Filterung und Sprachausgabe. Zusätzlich werden die Rahmenbedingungen beschrieben in welchen dieses Projekt zu verlaufen hat, um das Fokusprojekt erfolgreich zu absolvieren.
Die Bestandteile des Projektes sind Untersuchung der vorhandenen Hard- und Software, die Optimierung der Firmware und Software, die Geräuschunterdrückung, Sprachtranskription, messtechnisch Untersuchung und Entwicklung einer Prototypsoftwarelösung. 
Folgende Aufgaben sind klar abgrenzbar und werden nicht Bestandteil des Projektes sein:
Entwicklung neuer Hardware und eigener Sensoren, Vollständige Eliminierung sämtlicher Geräusche, Übersetzung in eine andere Sprache, Industrielle Produktzulassung und Zertifizierung als Gehörschutz.
Unser Projekt wird sich lediglich mit dem Forschen und Messen den oben genannten Zielen auseinandersetzen.


\pagebreak

\section{Projektorganisation}
Die Gruppe, bestehend aus sieben Studierenden, setzt sich aus folgenden Aufgaben und Untergruppen zusammen:
Die Aufgaben Projektmanager und Softwareimplementierung übernimmt Elmo Hofmann. Die Marktforschung und Erhebung wichtiger Fallzahlen übernimmt Tobias Holbach. Timon Kramer wird sich mit der Hardware der Geräte auseinandersetzten.


\subsection{Untergruppe: Speech-To-Text}
Die Untergruppe, deren Fokus auf Speech-to-Text liegt, setzt sich zusammen aus: \\
Die Hauptaufgaben dieser setzten sich zusammen aus:
\begin{itemize}
    \item Überprüfung der Forschungsergebnisse der vorherigen Gruppe
    \item Analyse der physikalischen Gegebenheit sowohl von sprachlicher- als auch textlicher Seite
    \item Erfassung der konkreten nötigen Softwareumgebung und Filterungsvorraussetzng 
    \item Zusammenarbeit mit der anderen Untergruppe in Bezug auf Signalverarbeitung und Filterung 
\end{itemize}

\subsection{Untergruppe: ANC}
Die Untergruppe, deren Fokus auf der Geräuschfilterung liegt, setzt sich zusammen aus Felix, Elmo, 
Die Hauptaufgaben dieser setzt sich zusammen aus:
\begin{itemize}
    \item Erfassung der technischen Gegebenheiten des Gerätes und notwendigen Signalfilterung
    \item Abgleichen des resultierenden möglichen Rahmens zur Umsetzung der Filtrung
    \item Erfassung der bisherigen Softwareimplementationen und die folgenden zuoptmierenden Stellen
    \item Absprache zuhalten mit der anderen Untergruppe in Bezug auf Signalverarbeitung und Filterung
\end{itemize}

\subsection{Gemeinsamen Aufgaben}
Die Aufgaben, die die ganze Gruppe betreffen, sind eine organisierte Arbeit der Themen und guter Kommunikation des Fortschrittes und geplanter Arbeitsschritte. Das Lastenheft, Pflichtenheft und die Dokumentation des Projektes mit zusätzlicher Präsentation wird von allen organisiert. 
Das große gemeinsame Ziel ist es die bestmögliche Filterung von Umgebungsgeräuschen im Bereich von Hz zu erfassen. Diesen Bereich wird im Rahmen der Software und technischen Rahmenbedingungen gilt es sop gut als möglich nicht hörbar zu machen für den Nutzer. Während dieser Bereich nun unterdrücktz wird, soll es nun möglich gemacht werden Gespräche zu verstärkend hörbar zu machen und diese an einer passenden Schnittstelle in Textform sichtbar zu machen.



\section{Anforderungen an das System}

\subsection{Allgemeine Anforderungen}
Die allgemeine Forderung, die wir an unsere Umsetzung setzten setzt sich wie folgt zusammen:
\begin{itemize}
    \item Die Unterdrückung der Geräusche im Bereich von Hz um ein Drittel zu steigern %todo richtigen Wert einfügen, das nur ein beispiel und keine gute beschreibung
    \item Die Textausgabe sicher von einfachen Gesprächen zu ermöglichen 
    \item Die Softwareimplementierung möglichst perfomant und laufzeitschonend umzusetzen
\end{itemize}

\subsection{Funktionale Anforderungen}
%todo mit den genauen werten was zb für geräusche und frequenzen gefiltert werden sollen, was für eine qualität der textausgabe und wie die software implementiert werden soll


\subsection{Nichtfunktionen Anforderungen}
%todo 
\begin{tabularx}{\textwidth}{ |>{\centering\arraybackslash}p{1.1cm}|>{\raggedright\arraybackslash}p{2.8cm}|>{\raggedright\arraybackslash}X| }
%\newcolumntype{s}{>{\columncolor[HTML]{AAACED}} p{3cm}}
\hline
\rowcolor{Turquoise} 
\arrayrulecolor[HTML]{DB5800}
%\multicolumn{3}{|c|}{Sensoren} \\
ID & Kategorie & Anforderung \\
\hline
         NF01 & Signalqualität & Verbesserung der Geräuschunterdrückung um mindestens X dB im definierten Frequenzbereich \\   %todo Wert einfügen
\hline
         NF02 & Latenz & Maximal zulässige Verzögerung der Audioverarbeitung\\
\hline
         NF03 & Genauigkeit & Maximal zulässige Wortfehlerrate von x Prozent der Transskription \\     % todo Wert einfügen
\hline
         NF04 & Ressourcen & Einhaltung der verfügbaren Rechen- und Speicherkapazitäten \\
\hline
         NF05 & Energieeffizienz & Einhaltung einer definerten Mindestbetriebdauer von x Stunden \\  %todo am bestsen einen Wert einfügen der für anwendung auf bau odedr  so sinnvoll ist
\hline
         NF06 & Sicherheit & Vermeidung gesundheitsgefährdender Schallpegelexposition \\
\hline
         NF07 & Datenschutz & Schutz personenbezogener Sprachaufzeichnungen \\
\hline
         NF08 & Wartbarkeit & Nachvollziehbarer und dokumentierter Quellcode\\
\hline
         NF09 & Zuverlässigkeit & Stabiler Betrieb während definierter Tagesräume \\
\hline
\end{tabularx}

\subsection{Gruppenspezifische Anforderungen}
Unsere Gruppe setzt sich eine möglichst effiziente und geordnete Arbeit in Untergruppen mit guter Kommunikation als wichtiges Ziel anzusehen. Jegliche Fortschritte und Ergebnisse werden in regelmäßig durchgeführten Sitzungen ausgetauscht und intermittend für alle Mitglieder einsehbar digital festgehalten.
Nach der Erhebung der theoretischen Erfassung aller nötigen Informationen werden diese kummuliert festgehalten und als Basis genutzt für die folgende Arbeit. 
Der Projektleiter stellt sicher, dass der Kommunikationskanal gepflegt wird und intergruppunell aufrecht gehalten wird. Zusätzlich setzt er sich für ein gutes Arbeisklima und das Wohlergehen aller Mitglieder ein. Jeder der Gruppenmitglieder halten sich an die Richtlinien der Nettiquette der Universität Ulm und dem respektvollen Umgang untereinander und den Ihnen zur Verfügung gestellten Materialien. Dies Beinhaltet die Unversehrtheit der Kopfhörer als auch des Datenschutzes.
Das wissenschaftliche Arbeiten, Erlernen von neuem Wissen und das Aufrechterhaltens eines guten wissenschalftlichen Forschungsstandards steht an erster Stelle.


\subsection{Priorisierung der Anforderungen}
Als oberstes Ziel steht das Gruppenverständnis, dass Gemeinsam an einem großen Ziel arbeitet wird, selbst wenn die Aufgaben der einzelnen Mitglieder sich auch vonaneinader unterscheiden. Eine gute Kommunikation und ständige Unterrichtung Zwischenstände sind als verständlich anzusehen. 
Danach folgt die Arbeit an Signalverarbeitung und -filterung mit anschlißenender Textausgabe.
\begin{tabularx}{\textwidth}{ |>{\centering\arraybackslash}p{1.1cm}|>{\raggedright\arraybackslash}X|>{\centering\arraybackslash}p{1.7cm}| }
%\newcolumntype{s}{>{\columncolor[HTML]{AAACED}} p{3cm}}
\hline
\rowcolor{Turquoise} 
\arrayrulecolor[HTML]{DB5800}
%\multicolumn{3}{|c|}{Sensoren} \\
ID & Anforderung & Priorität \\
\hline
         F01 & System muss akustische Umgebungssignale erfassen können & Muss \\
\hline
         F02 & System muss störende Umgebungsgeräusche in einem definierten Frequenzbereich reduzieren & Muss \\
\hline
         F03 & System muss gesprochene Sprache erfassen und in Text umwandeln können & Muss \\
\hline
         F04 & System muss den erzeugten Text auf einem externen Endgerät darstellen können mit einer gewissen Qualität & Muss \\
\hline
         F05 & System muss relevante Messdaten digital aufzeichnen können & Muss \\
\hline
          F06 & Sprachverständlichkeit soll bei vorhandenen Störgeräuschen verbessert werden & Soll \\
\hline
          F07 & Die Geräuschunterdruckung soll in ihrer Intensität anpassbar sein & Soll \\
\hline
         F08 & System soll in der Lage sein die restliche Belastung der Geräusche zu bewerten und eine Warnung bei einem Grenzwert zu auslösen & Kann \\
\hline
\end{tabularx}

\section{Technische Rahmenbedingungen}

\subsection{Hardware}
Vom Hersteller sind folgende Bestandteile, neben der Aufbewahrung und austauschbaren Ohrpads, in den Kopfhörern verbaut:

\begin{itemize}
    \item Folgende Bestandteile sind Sensoren:
        \subitem  Eine 9-Achsen Bewegungssensor (zusätzliche Messung möglicher Deformationen des Ohrkanals)
        \subitem ein Drei-Wellenlänge Photoplethysmograph (infratrot, rot und grün PPG Sensoren + Pulsoximeter)
        \subitem einen Drucksensor (Messung des Innendruck des Hörkanals) 
        \subitem Zwei Mikrophone (Luftübertragung und Knochenübertragung, Ultra- und Infrasound, Schnittstelle zur Aufnahme oder des Streamings)
        \subitem einen optischen Temperatursensor (Messung Hauttemperatur)
    \item Restliche verbaute Bestandteile:
        \subitem Eine Knopfbatterie
        \subitem Einen USB-Anschluss
        \subitem einen Druckknopf und eine RGB-LED
        \subitem Dynamic Driver-Lautsprecher
    \item Das Außengehäuse
\end{itemize}  
    
\subsection{Software}
Die Software wird vom GitHub des Herstellers bezogen %todo einfügen des links als quelle vom github

\subsection{Schnittstelle}
Die Kopfhörer sind sowohhl Speicher für die Software als auch Träger der Sensoren, somit wird die Aufzeichnung der Messungen und Daten lediglich auf einer extern engefügten SD-Karte gespeichert falls geofordert oder wird kabellos gestreamed.

\subsection{Messgrößen und Messbereiche}

\begin{tabularx}{\textwidth}{ |>{\raggedright\arraybackslash}p{4.1cm}|>{\centering\arraybackslash}p{3.2cm}|>{\centering\arraybackslash}X| }
%\newcolumntype{s}{>{\columncolor[HTML]{AAACED}} p{3cm}}
\hline
\rowcolor{Thistle} 
\arrayrulecolor[HTML]{DB5800}
%\multicolumn{3}{|c|}{Sensoren} \\
Sensor & Micro-SD Aufzeichnung & Kabellose Übertragung \\
\hline
         Luftmikrophon (Inner-/Außerohrs) & 48 kHz & 48 kHz \\
\hline
         Knochenübertragungssensor der Beschleunigungen & 6,4 kHz & 800 Hz \\
\hline
         PPG-Sensor & 4,096 kHz & 512 Hz \\
\hline
         9-Achsen Bewegungsensor & 800 Hz & 100 Hz \\
\hline
         Temperatursensor & 64 Hz & 64 Hz \\
\hline
         Ohrkanaldrucksensor & 200 Hz & 200 Hz \\
\hline
\end{tabularx}

\section{Validierung und Abnahmekriterien}

    \subsection{Mess- und Testbedingungen}
    Die Messungen werden nach einem standatisiertem Schema ablaufen. Der Ort wird in einem Raum mit definierten akustischen Eigenschaften stattfinden. Die Messungen werden mit einer definierten Anzahl an Testpersonen durchgeführt, die in einem definierten Abstand zu den Lautsprechern sitzen. Die Lautsprecher werden definierte Geräusche und Sprache in einem definierten Frequenzbereich abspielen. Die Messungen werden sowohl mit als auch ohne der Kopfhörer durchgeführt, um die Wirksamkeit der Geräuschunterdrückung und Sprachtranskription zu bewerten. 
    Die Messgeräte sind geeichte und kalibrierte Geräte, die die Messungen der Geräuschpegel und der Sprachverständlichkeit durchführen. Die Messungen werden in einem definierten Zeitraum durchgeführt, um die Konsistenz der Ergebnisse zu gewährleisten. Die Testpersonen werden vor den Messungen über die Vorgehensweise informiert und ihre Zustimmung zur Teilnahme eingeholt. Die Messungen werden unter Einhaltung der geltenden Datenschutzbestimmungen durchgeführt, um die Privatsphäre der Testpersonen zu schützen.
    Zudem werden die Messungen unter verschiedenen Bedingungen durchgeführt, um die Robustheit der Geräuschunterdrückung und Sprachtranskription zu bewerten. Dazu gehören unterschiedliche Lautstärken, Frequenzen und Arten von Hintergrundgeräuschen. Die Ergebnisse der Messungen werden dokumentiert und analysiert, um die Wirksamkeit der Geräuschunterdrückung und Sprachtranskription zu bewerten und mögliche Verbesserungen zu identifizieren.
    Jegliche Messungen werden in einem standardisierten Format dokumentiert, um die Vergleichbarkeit der Ergebnisse zu gewährleisten. Die Ergebnisse werden in einem Abschlussbericht zusammengefasst, der die Wirksamkeit der Geräuschunterdrückung und Sprachtranskription bewertet und Empfehlungen für zukünftige Verbesserungen enthält.
    Messunsicherheiten sind minimiert, jedoch nicht vollständig ausgeschlossen. Sowohl die Messgeräte als auch die Testpersonen können zu Abweichungen in den Ergebnissen führen. Die Messungen werden daher mehrfach durchgeführt, um die Konsistenz der Ergebnisse zu gewährleisten und mögliche Ausreißer zu identifizieren. Die Ergebnisse werden statistisch analysiert, um die Signifikanz der Unterschiede zwischen den Messungen mit und ohne Kopfhörer zu bewerten. 
    Die einzelnen Komponenten der Kopfhörer haben eine definierte Toleranz, die in den Messungen berücksichtigt wird. Die Ergebnisse der Messungen werden in einem standardisierten Format dokumentiert, um die Vergleichbarkeit der Ergebnisse zu gewährleisten.
    In dem Dokument, das alle Messungen und Ergebnisse zusammenfasst, werden die Messbedingungen, Testpersonen, Messgeräte, Messunsicherheiten und Toleranzen der Kopfhörerkomponenten detailliert beschrieben. Die Ergebnisse werden in Tabellen und Diagrammen dargestellt, um die Wirksamkeit der Geräuschunterdrückung und Sprachtranskription zu bewerten und mögliche Verbesserungen zu identifizieren. Der Abschlussbericht wird den Auftraggebern zur Verfügung gestellt, um die Ergebnisse der Messungen und die Wirksamkeit der Geräuschunterdrückung und Sprachtranskription zu bewerten.

    \subsection{Bewertung der Sprachtranskription}
    Die Bewertung der Sprachtranskription erfolgt anhand der Wortfehlerrate (WER), die den Prozentsatz der falsch erkannten Wörter im Vergleich zur Gesamtzahl der gesprochenen Wörter angibt. Eine niedrige WER deutet auf eine hohe Genauigkeit der Transkription hin. Die WER wird durch den Vergleich der transkribierten Texte mit den Originaltexten berechnet, wobei falsch erkannte, ausgelassene und hinzugefügte Wörter berücksichtigt werden. Die Bewertung erfolgt unter verschiedenen Bedingungen, einschließlich unterschiedlicher Lautstärken, Frequenzen und Arten von Hintergrundgeräuschen, um die Robustheit der Sprachtranskription zu bewerten. Die Ergebnisse werden in einem standardisierten Format dokumentiert und analysiert, um die Wirksamkeit der Sprachtranskription zu bewerten und mögliche Verbesserungen zu identifizieren.
    Im Vorhinein wird ein festes Skript der zu sprechenden Wörter und Sätze erstellt, welches von den Testpersonen während der Messung vorgelesen wird. Dieser Text beinhaltet eine definierte Anzahl an Wörtern in Sätzen, die steigernd in ihrer Komplexität und Wortvarianz sind. Diese Sätze werden in einem definierten Abstand zu den Lautsprechern vorgelesen, die die Hintergrundgeräusche abspielen. Die Testpersonen werden vor den Messungen über die Vorgehensweise informiert und ihre Zustimmung zur Teilnahme eingeholt. Die Messungen werden unter Einhaltung der geltenden Datenschutzbestimmungen durchgeführt, um die Privatsphäre der Testpersonen zu schützen. 

    \subsection{Bewertung der Geräuschunterdrückung}
    Die Bewertung der Geräuschunterdrückung erfolgt durch die Messung des Schalldruckpegels (SPL) in Dezibel (dB) vor und nach der Anwendung der Geräuschunterdrückung. Eine effektive Geräuschunterdrückung sollte zu einer signifikanten Reduktion des SPL führen, insbesondere in den definierten Frequenzbereichen, die als störend identifiziert wurden. Die Messungen werden unter verschiedenen Bedingungen durchgeführt, einschließlich unterschiedlicher Lautstärken, Frequenzen und Arten von Hintergrundgeräuschen, um die Robustheit der Geräuschunterdrückung zu bewerten. Die Ergebnisse werden in einem standardisierten Format dokumentiert und analysiert, um die Wirksamkeit der Geräuschunterdrückung zu bewerten und mögliche Verbesserungen zu identifizieren.
    Diese Geräusche werden vorerst von einem Lautsprecher abgespielt, der in einem definierten Abstand zu den Testpersonen positioniert ist. Die Testpersonen werden vor den Messungen über die Vorgehensweise informiert und ihre Zustimmung zur Teilnahme eingeholt. Die Messungen werden unter Einhaltung der geltenden Datenschutzbestimmungen durchgeführt, um die Privatsphäre der Testpersonen zu schützen. Jegliche Messungen werden in einem standardisierten Format dokumentiert, um die Vergleichbarkeit der Ergebnisse zu gewährleisten. Die Ergebnisse werden in Tabellen und Diagrammen dargestellt, um die Wirksamkeit der Geräuschunterdrückung zu bewerten und mögliche Verbesserungen zu identifizieren.
    Die Messdaten wird auf einer externen SD-Karte gespeichert und sind somit standatisiert und für eine Auswertung längerfristig verfügbar.


\section{Zeitplan und Meilensteine}
5.11.2026 Kick-Off Meeting\\
26.11.2026 Besprechung der festgelegten Ziele und festlegen des neuen Termins für die Besprechung der Aufgabenverteilung und der Untergruppenaufteilung\\
10.12.2026 Besprechung der Aufgabenverteilung und der Untergruppenaufteilung und Festlegung der nächsten Schritte\\


\section{Risikoeinschätzung}
Jegliche Risiken, die durch die Nutzung der Kopfhörer entstehen können, sind von dem Hersteller vordefiniert und werden in den Kopfhörern durch die Software und Hardware des Herstellers geregelt. Jegliche Risiken, die durch die Nutzung der Software entstehen können, sind von uns zu definieren und zu bewerten. 

    \subsection{Gesundheitliche Risiken}
%todo 

    \subsection{Datenschutz}
    Alle Daten, die durch die Nutzung der Kopfhörer, auch im Rahmen der Messungen, entstehen, werden auf einer externen SD-Karte gespeichert. 
    Diese Daten werden nicht an Dritte weitergeben und auch nicht im Internet veröffentlicht.
    Alle Teilnehmer der Messungen werden annonymisiert und es werden keine personenbezogenen Daten gespeichert. 
    Jegliche Daten, die im Rahmen der Messungen entstehen, werden nur für die Auswertung der Messungen genutzt und nach Abschluss des Projektes gelöscht.
    Die Teilnehmer der Messungen werden über die Vorgehensweise informiert und ihre Zustimmung zur Teilnahme eingeholt. 
    Die Messungen werden unter Einhaltung der geltenden Datenschutzbestimmungen durchgeführt, um die Privatsphäre der Testpersonen zu schützen.
    Jegliche Daten, die vom Hersteller auf den Kopfhörer sind für uns nicht einsehbar.
    Diese PRojektarbeit wird unter Einhaltung der geltenden Datenschutzbestimmungen der Universität Ulm und der DSGVO durchgeführt. Jegliche Daten der Messungen und personenbezogene Daten werden nach Abschluss des Projektes werden bei allen Teilnehmern gelöscht und lediglich die Auftraggeber besitzen diese Daten.

    \subsection{Projektbezogene Risiken}
    Die Projektbezogenen Risiken beschränken sich nur während dem Tragen der Geräte der Messungen. Diese können sein: Verletzungen im Hörkanal einschließlich des Trommelfelles oder der Hörsinnzellen, Hörverlust(temporär oder permanent), Kopfschmerzen, Unwohlsein, Übelkeit, Schwindel, Müdigkeit, Konzentrationsschwierigkeiten, Orientierungsverlust, Angst, Hautausschlag, Allergische Reaktionen.
    Alle Versuche werden so durchgeführt, dass die Risiken minimniert werden. Das beinhaltet auch das Exponieren der lauten Geräusche auf einem Minimum zu halten und keine gesundheitsschädinge Grenze zu überschreiten. 
    Versuchsteilnehmer werden regelmäßige Pausen zur Verfügung gestellt und die Dauer der Messungen wird auf ein Minimum reduziert. Jegliche Risiken werden vor den Messungen den Teilnehmern erläutert und ihre Zustimmung zur Teilnahme eingeholt. Die Messungen werden unter Einhaltung der geltenden Sicherheitsbestimmungen durchgeführt, um die Gesundheit der Testpersonen zu schützen. 



\section{Erwartete Projektergebnisse und Deliverables}
Die erwarteten Projektergebnisse sind, dass die Geräuschunterdrückung im definierten Frequenzbereich um mindestens X dB verbessert wird, die Wortfehlerrate der Sprachtranskription unter Y Prozent liegt und die Software implementiert ist, die die Geräuschunterdrückung und Sprachtranskription in Echtzeit ermöglicht.
Die Auftrageber werden von uns ein Dokument erhalten, das die Ergebnisse alles Messungen und Arbeitsschritte der Gruppenmitglieder beinhaltet. 


\end{document}